%!TEX program = xelatex
\PassOptionsToPackage{full}{textcomp}
\documentclass[lang=en,11pt,a4paper,citestyle=authoryear,nobib=true]{elegantpaper}

\title{HW7 - STAT 709}
\author{Wells Guan}

\usepackage{array,url,stix}
\usepackage{tikz}
\usepackage{tikz-cd}
\usepackage{enumitem}
\usepackage{amsmath}
\usepackage{amssymb}
\usepackage{listings}
\usepackage{subfigure}

\lstset{
    frame=single,
    columns=fullflexible,
    keepspaces=true,
    showstringspaces=false,
    breaklines=true,
    breakatwhitespace=false,
    basicstyle=\ttfamily\footnotesize,
    numbers=left,
    numberstyle=\tiny,
    numbersep=6pt
}

\newcommand{\Z}{\mathbb{Z}}
\newcommand{\R}{\mathbb{R}}
\newcommand{\N}{\mathbb{N}}
\newcommand{\B}{\mathcal{B}}
\newcommand{\E}{\mathbb{E}}
\renewcommand{\P}{\mathbb{P}}
\newcommand{\norm}[1]{\left\lVert#1\right\rVert}

\begin{document}
\maketitle

\noindent All four exercises are required. Complete every part of each exercise.

\subsection*{Problem 1. Modes of convergence}
Jun Shao, \textit{Mathematical Statistics}, second edition, Chapter 1:
Exercises 116 and 141 (required).

\subsubsection*{Exercise 1.116}
Let $X,X_1,X_2,\ldots$ be random variables. Find an example for each of
the following cases:
\begin{enumerate}[label=(\alph*),leftmargin=*]
\item $X_n\xrightarrow{p}X$, but $\{X_n\}$ does not converge to $X$ a.s.
\item $X_n\xrightarrow{p}X$, but $\{X_n\}$ does not converge to $X$ in $L^p$
for any $p>0$.
\item $X_n\xrightarrow{d}X$, but $\{X_n\}$ does not converge to $X$ in
probability (do not use Example 1.26).
\item $X_n\xrightarrow{p}X$, but $\{g(X_n)\}$ does not converge to $g(X)$ in
probability for some function $g$.
\end{enumerate}

\par\vspace{0.5em}\noindent\textbf{Sol.}\par
For all the $X,X_i$, we assume them to be the measurable function from $([0,1],\mathcal{B}([0,1])) \to (\R,\mathcal{B})$.
\begin{enumerate}[label=(\alph*),leftmargin=*]
\item Consider \[X_n = \chi\left(\left(\dfrac{n-2^{k(n)}}{2^{k(n)}},\dfrac{n+1-2^{k(n)}}{2^{k(n)}}\right]\right)\] where $k(n) = \left[\log_2 n\right]$ with $X = 0$.
\item Consider
\[
X_n = n\chi((0,1/(\ln n)])
\]
with $X = 0$.
\item Consider
\[
X_n = (1-x)
\]
with $X = x$.
\item Consider
\[
X_n = 1/n
\]
and $X = 0$ with $g(x) = \chi((0,\infty))$.
\end{enumerate}

\subsubsection*{Exercise 1.141}
Suppose that $X_n$ is a random variable having the binomial distribution
$Bi(p_n,n)$. Show that if $np_n\to\theta>0$, then $X_n\xrightarrow{d}X$,
where $X$ has the Poisson distribution $P(\theta)$.

\par\vspace{0.5em}\noindent\textbf{Sol.}\par\vspace{3em}
Compute the distribution function of $X_n$, which is
\[
\begin{aligned}
F_{X_n}(x) &= \sum\limits_{m=0}^{[x]} C_n^mp_n^{m}(1-p_n)^{n-m} \\
&= \sum\limits_{m=0}^{[x]} \dfrac{(np_n)^m}{m!}(1-p_n)^{n}(n!/(n-m)!n^m)(1-p_n)^{-m} \\
&\to \sum\limits_{m=0}^{[x]} \dfrac{\theta^me^{-\theta}}{m!} = F_x(x),\quad n\to\infty
\end{aligned}
\]
for any $x \in \R$ for $X \sim P(\theta)$ since
\[
n\ln(1-p_n) \sim -np_n \to -\theta.
\]
\newpage
\subsection*{Problem 2. Asymptotic distributions and estimation}
Bickel and Doksum, \textit{Mathematical Statistics: Basic Ideas and Selected
Topics}, Volume I, second edition, Chapter 5, Section 5.6,
Problems for Section 5.3: Exercises 25 and 26 (required).

\subsubsection*{Exercise 5.3.25}
Let $X_1,\ldots,X_n$ be a sample from a population with
$\sigma^2=\operatorname{Var}(X)<\infty$, $\mu=E(X)$ and let
$T=\bar X^2$ be an estimate of $\mu^2$.
\begin{enumerate}[label=(\alph*),leftmargin=*]
\item When $\mu\neq0$, find the asymptotic distribution of
$\sqrt n(T-\mu^2)$ using the delta method.
\item When $\mu=0$, find the asymptotic distribution of $nT$ using
\[
P(nT\leq t)=P\bigl(-\sqrt t\leq\sqrt n\bar X\leq\sqrt t\bigr).
\]
Compare your answer to the answer in part (a).
\item Find the limiting laws of $\sqrt n(\bar X-\mu)^2$ and
$n(\bar X-\mu)^2$.
\end{enumerate}

\par\vspace{0.5em}\noindent\textbf{Sol.}\par
\begin{enumerate}[label=(\alph*),leftmargin=*]
\item By the CLT, we know
\[
\sqrt{n}(\overline{X} - \mu) \Rightarrow \mathcal{N}(0,\sigma^2)
\]
and hence by Delta method, we have
\[
\sqrt{n}(T - \mu^2) \Rightarrow \mathcal{N}(0,4\mu^2\sigma^2).
\]
\item We know that
\[
P(\sqrt{n}T \leq t) = P(nT \leq \sqrt{n}t) = P(-\sqrt{\sqrt{n}t} \leq \sqrt{n}\overline{X} \leq \sqrt{\sqrt{n}t}) \to 1
\]
with $n \to \infty$ and hence $\sqrt{n}T \Rightarrow 0$ since $T$ is non-negative, which is coinciding the conclusion in (a) by setting $\mu = 0$.
\item Assume $Y_i = X_i - \mu$ and then we know $\overline{X} - \mu = \overline{Y}$. We may use (a) for $T' = \overline{Y}$ and obtain that
\[
\sqrt{n}(\overline{X} - \mu)^2 \Rightarrow 0.
\]
Then we use the continuous mapping and hence
\[
{n}(\overline{X} - \mu)^2 \Rightarrow \sigma^2\chi_1
\]
where $\chi_1$ has the standard Chi-square distribution.
\end{enumerate}

\subsubsection*{Exercise 5.3.26}
Show that if $X_1,\ldots,X_n$ are i.i.d.\ $\mathcal N(\mu,\sigma^2)$, then
\[
\sqrt n\bigl(\bar X-\mu,\widehat\sigma^2-\sigma^2\bigr)
\xrightarrow{\mathcal L}\mathcal N(0,0,\Sigma_0)
\]
where $\Sigma_0=\operatorname{diag}(\sigma^2,2\sigma^4)$.

\noindent\textit{Hint:} Use (5.3.33) and Theorem 5.3.4.

\par\vspace{0.5em}\noindent\textbf{Sol.}\par\vspace{3em}
Notice that
\[
\hat{\sigma}^2 = \frac{1}{n}\sum_{i=1}^n (X_i-\bar X)^2.
\]
Expanding around $\mu$ gives
\[
\hat{\sigma}^2
= \frac{1}{n}\sum_{i=1}^n (X_i-\mu)^2 - (\bar X-\mu)^2.
\]
Thus,
\[
\sqrt n
\begin{pmatrix}
\bar X-\mu\\
\hat{\sigma}^2-\sigma^2
\end{pmatrix}
=
\frac{1}{\sqrt n}\sum_{i=1}^n
\begin{pmatrix}
X_i-\mu\\
(X_i-\mu)^2-\sigma^2
\end{pmatrix}
-
\begin{pmatrix}
0\\
\sqrt n(\bar X-\mu)^2
\end{pmatrix}.
\]
By the CLT, $\sqrt n(\bar X-\mu)=O_p(1)$, so
\[
\sqrt n(\bar X-\mu)^2
= \frac{\bigl(\sqrt n(\bar X-\mu)\bigr)^2}{\sqrt n}
\xrightarrow{p} 0.
\]

The summands are i.i.d.\ random vectors with mean zero.
To compute their covariance matrix, let $Y\sim\mathcal N(0,\sigma^2)$.
Since $\mathbb E[Y^3]=0$ and $\mathbb E[Y^4]=3\sigma^4$, we have
\[
\begin{aligned}
\operatorname{Var}(Y) &= \sigma^2,\\
\operatorname{Cov}(Y,Y^2-\sigma^2) &= \mathbb E[Y^3] = 0,\\
\operatorname{Var}(Y^2-\sigma^2)
&= \mathbb E[Y^4] - \bigl(\mathbb E[Y^2]\bigr)^2
= 2\sigma^4.
\end{aligned}
\]
Therefore, by the CLT and Slutsky's theorem,
\[
\sqrt n
\begin{pmatrix}
\bar X-\mu\\
\hat{\sigma}^2-\sigma^2
\end{pmatrix}
\xrightarrow{\mathcal L}
\mathcal N_2\!\left(
\begin{pmatrix}0\\0\end{pmatrix},
\begin{pmatrix}
\sigma^2 & 0\\
0 & 2\sigma^4
\end{pmatrix}
\right)
\]
we are done.
\end{document}
